\documentclass[11pt,a4paper]{article}

\usepackage{amsmath,amssymb}
\usepackage{booktabs}
\usepackage{hyperref}
\usepackage{geometry}
\usepackage{microtype}
\usepackage{parskip}
\usepackage{listings}
\usepackage{xcolor}

\geometry{margin=2.5cm}

\lstset{
  basicstyle=\ttfamily\small,
  backgroundcolor=\color{gray!8},
  frame=single,
  framerule=0.4pt,
  breaklines=true,
  language=Python,
}

\title{\textbf{vzlab: A Four-Tier Multi-Species Neural Avatar\\Validation Framework}\\[6pt]
  \large Technical Note}
\author{vzebra project}
\date{2026-04-25}

\begin{document}
\maketitle

\begin{abstract}
We describe \textsc{vzlab}, an open Python framework for building and
validating virtual-organism neural avatars across species.
The framework enforces a four-tier validation protocol: behavioural
battery~(T1), atlas rank correspondence~(T2), lesion replication~(T3),
and sensory dropout robustness~(T4).
All four current species implementations---\textit{Danio rerio}
(7,316 spiking neurons), \textit{Caenorhabditis elegans}~(302),
\textit{Drosophila melanogaster}~(2,145), and \textit{Xenopus laevis}
tadpole~(150)---achieve Grade~A+ (all four tiers pass), spanning three
phyla and four orders of magnitude in neuron count.
Robustness indices under 50\,\% sensory dropout range from 0.946
(\textit{Drosophila}) to 1.000 (zebrafish and xenopus), confirming
fault-tolerant distributed representations.
\end{abstract}

%----------------------------------------------------------------------
\section{Introduction}

Computational neuroscience models are typically validated against
single experiments in a single species.
\textsc{vzlab} introduces a systematic cross-species validation
paradigm in which the same four-tier protocol is applied to every
neural avatar, enabling direct comparison of model fidelity across
organisms.

The framework was designed around three principles:
\begin{enumerate}
  \item \textbf{Species-agnostic contracts.}  All species share four
    abstract base classes so that experimental protocols, lesion tools,
    and analysis pipelines require no modification when the species
    changes.
  \item \textbf{Graded, reproducible validation.}  Tier grades
    (A+/A/B) encode what kind of biological evidence the avatar
    reproduces and are serialised with every result object.
  \item \textbf{Fault-tolerant representations.}  Tier~4 tests whether
    the avatar's behaviour degrades gracefully under sensory loss,
    matching the redundancy of biological neural circuits.
\end{enumerate}

%----------------------------------------------------------------------
\section{Architecture}

\subsection{Abstract Base Classes}

Four ABCs define the species-agnostic platform contracts
(\texttt{vzlab/core/interfaces.py}):

\begin{description}
  \item[\texttt{Connectome}] Wraps anatomical data---region names,
    neuron counts, projection matrices, neuromodulatory targets.
    Ingests FlyWire, ZFIN, MICrONS, and the White~\textit{et~al.}\
    \textit{C.\ elegans} connectome.
  \item[\texttt{BrainModule}] Implements one timestep of neural
    dynamics (\texttt{step()}), returns a six-level
    \texttt{HierarchicalState}, and exposes
    \texttt{ablate()}/\texttt{restore()}/\texttt{set\_param()}.
  \item[\texttt{Environment}] Emits typed sensory structs:
    \texttt{PhotonField}, \texttt{ChemField}, \texttt{PressureField},
    \texttt{BodyState}, \texttt{SocialField}.
  \item[\texttt{BehavioralProtocol}] Defines stimuli, episode length,
    and a \texttt{score()} method converting world state to metrics.
\end{description}

\subsection{ExperimentRunner}

The \texttt{ExperimentRunner} composes a \texttt{VirtualOrganism}
(brain + environment) with a protocol and exposes a chainable
intervention API:

\begin{lstlisting}
from vzlab import build
from vzlab.engine.experiment import ExperimentRunner
from vzlab.protocols import PredatorAvoidanceProtocol

organism = build("danio_rerio")
runner = ExperimentRunner(organism, PredatorAvoidanceProtocol(), seed=42)
runner.ablate("amygdala")
result = runner.run(n_episodes=20)
print(result.summary())
\end{lstlisting}

The identical API works for all four species.

\subsection{Environments}

\begin{table}[h]
\centering
\caption{Environments per species.}
\begin{tabular}{lll}
\toprule
Species & Environment & Key parameters \\
\midrule
\textit{D.\ rerio}         & \texttt{AquaticArena}  & 800$\times$600\,px, 12 food, 1 predator \\
\textit{C.\ elegans}       & \texttt{AgarPlate}     & 100$\times$100\,mm, Gaussian chem.\ gradients \\
\textit{D.\ melanogaster}  & \texttt{OdourArena}    & 72-odour CS+/CS$-$ conditioning \\
\textit{X.\ laevis}        & \texttt{SwimTank}      & 400$\times$300\,px channel, wall-touch CPG \\
\bottomrule
\end{tabular}
\end{table}

%----------------------------------------------------------------------
\section{Four-Tier Validation Protocol}

\subsection{Tier 1 — Behavioural Battery}

Each species has a curated set of behavioural tests run via
\texttt{Tier1Battery.run(species)}.
Tests check survival competence, species-characteristic motor
patterns, and sensory-motor coupling against published thresholds.

Zebrafish tests include: prey capture kinematics (J-turn, capture
strike), predator-charge escape, obstacle avoidance, shoaling
cohesion (3~agents, cohesion $\geq 0.20$), and social alarm
propagation (alarm propagation fraction $\geq 0.50$).

\subsection{Tier 2 — Atlas Rank Correspondence}

Regional activity rates returned in \texttt{HierarchicalState.L4\_region}
are rank-correlated (Spearman's $r_s$) against species-specific
biological reference vectors derived from whole-brain imaging or
electrophysiology atlases.
Pass criterion: $r_s > 0.50$ ($p < 0.05$).

\subsection{Tier 3 — Lesion Replication}

At least one published lesion experiment is replicated per species.
The avatar must reproduce the sign and approximate magnitude of the
reported deficit.

\begin{table}[h]
\centering
\caption{Tier~3 lesion tests.}
\begin{tabular}{lll}
\toprule
Species & Lesion & Expected deficit \\
\midrule
\textit{D.\ rerio}         & Habenula ablation          & Reduced goal-switching ($\geq$20\%) \\
\textit{C.\ elegans}       & AWC ablation               & CI drops from 0.68 to $<$0.10 \\
\textit{D.\ melanogaster}  & DAN silencing              & Discrimination drops from 0.72 to $<$0.30 \\
\textit{X.\ laevis}        & dIN ablation               & Swimming speed drops $\geq$80\% \\
\textit{X.\ laevis}        & cIN ablation               & Alternation fraction drops $\geq$0.30 \\
\bottomrule
\end{tabular}
\end{table}

\subsection{Tier 4 — Sensory Dropout Robustness}

Bernoulli dropout is applied to the sensory signal struct
\emph{before} \texttt{brain.step()}, so the avatar makes decisions
from genuinely degraded input.
Dropout fractions $p \in \{0.10, 0.30, 0.50\}$ are tested over 5
seeds $\times$ 40 steps each.
The retention ratio at fraction $p$ is:
\[
  \rho(p) = \frac{\bar{S}(p)}{\bar{S}(0)},
\]
where $\bar{S}$ is the mean behavioural score.
Pass thresholds: $\rho(0.10) \geq 0.80$, $\rho(0.30) \geq 0.50$,
$\rho(0.50) \geq 0.20$.

The Robustness Index aggregates performance across all dropout levels:
\[
  \text{RI} = \frac{\displaystyle\int_0^{0.5} \rho(p)\,dp}{0.5}
  \approx \frac{\text{trapezoid}(\rho,\, [0,0.1,0.3,0.5])}{0.5}.
\]
$\text{RI} = 1$ means perfect retention at all dropout levels;
$\text{RI} = 0$ means complete failure at $p = 0$.

\paragraph{Implementation (\texttt{vzlab/validation/tier4.py})}

\begin{lstlisting}
def _apply_sensory_dropout(signals, p, rng):
    s = deepcopy(signals)
    if s.chem:
        mask = rng.random(len(s.chem.concentrations)) >= p
        s.chem.concentrations *= mask
    if s.body and rng.random() < p:
        s.body.in_contact = False
    if s.photon:
        mask = rng.random(s.photon.intensity.shape) >= p
        s.photon.intensity *= mask
    if s.pressure:
        mask = rng.random(len(s.pressure.flow_vectors)) >= p
        s.pressure.flow_vectors *= mask[:, None]
    return s
\end{lstlisting}

\subsection{Grading}

\begin{description}
  \item[Grade A+] T1 + T2 + T3 + T4 all pass.
  \item[Grade A]  T1 + T2 + T3 pass; T4 not run or partial.
  \item[Grade B]  T1 + T2 pass.
  \item[Grade C]  T1 passes only.
\end{description}

%----------------------------------------------------------------------
\section{Species Implementations}

\subsection{\textit{Danio rerio} (larval zebrafish)}

\textbf{Architecture.}
7,316 Izhikevich spiking neurons across 48 brain modules (L1--L6).
Seven cell types (RS, IB, FS, CH, LTS, TC, MSN); 75\,\%~E / 25\,\%~I.
Neuron positions constrained by the MPIN whole-brain atlas
(47,010~cells, 72~bilateral regions).
Four-axis neuromodulation (DA, NA, 5-HT, ACh).

\textbf{Goal selection.}
Active inference Expected Free Energy minimisation across four
policies (forage, flee, explore, social) with reward-modulated
three-factor STDP eligibility traces ($\tau = 500$\,ms).

\textbf{STDP synaptic dropout.}
Each \texttt{EligibilitySTDP.consolidate()} call applies a Bernoulli
mask to the weight delta before committing (default $p_\text{drop} =
0.10$), preventing synapse co-adaptation and forcing redundant
distributed representations analogous to biological synaptic release
failure ($\sim$10--20\,\% per terminal).

\textbf{Validation.}
T1: 7/7 behavioural tests (including shoaling cohesion and social
alarm propagation).
T2: $r_s = 0.81$ (MPIN atlas).
T3: habenula lesion reduces goal-switching by 22\,\%.
T4: RI = 1.000.
\textbf{Grade~A+.}

\subsection{\textit{Caenorhabditis elegans}}

\textbf{Architecture.}
302 leaky integrate-and-fire neurons ($\tau = 20$\,ms,
$\theta = 1.0$), 7,391 chemical plus 890 electrical synapses from
the White~\textit{et~al.}\ connectome.

\textbf{Validation.}
T1: chemotaxis index CI $= 0.68$ (threshold $\geq 0.60$).
T2: $r_s = 0.72$ (activity rank vs.\ published profiles).
T3: AWC ablation drops CI from 0.68 to $0.04 \pm 0.02$ (biological:
$\leq 0.10$).
T4: RI = 0.960.
\textbf{Grade~A+.}

\subsection{\textit{Drosophila melanogaster}}

\textbf{Architecture.}
2,145 neurons: 50~PNs, 700~KC$_{\alpha\beta}$, 300~KC$_{\alpha'\beta'}$,
1,000~KC$_\gamma$, 1~APL (global WTA inhibition, 10\,\% sparsity),
60~DANs, 34~MBONs.
Connectivity from the FlyWire hemibrain connectome.

\textbf{Sparsity enforcement.}
APL feedback silences all KCs below the 90th-percentile firing
threshold each timestep, implemented by counting active units and
zeroing the bottom 90\,\% by rate.

\textbf{Validation.}
T1: odour discrimination $= 0.72$ (threshold $\geq 0.60$).
T2: $r_s = 0.68$ (MB calcium imaging reference).
T3: DAN silencing drops discrimination from 0.72 to
$0.19 \pm 0.06$ (threshold $\leq 0.30$).
T4: RI = 0.946.
\textbf{Grade~A+.}

\subsection{\textit{Xenopus laevis} tadpole}

\textbf{Architecture.}
150 neurons: 20~Rohon--B\'{e}ard (RB) sensory, 60~descending
interneurons (dIN, pacemakers), 40~commissural interneurons (cIN,
L/R alternation), 30~motoneurons (MN).
Circuit topology from Roberts~\textit{et~al.}\ 1998.

\textbf{CPG dynamics.}
Touch $\to$ RB $\to$ dIN rhythm (tonic firing rate 0.15--0.35
during steady swimming, not maximal at every wall contact).
cIN cross-inhibition enforces left--right alternation.
MN drives tail motor output.

\textbf{Validation.}
T1: swimming speed $\geq 0.50$, alternation fraction $\geq 0.80$.
T2: $r_s = 1.00$ (rb $<$ cIN $<$ MN $<$ dIN rank order).
T3: dIN ablation $\to$ 100\,\% speed reduction (threshold $\geq 80$\,\%);
cIN ablation $\to$ alternation drops by 0.45 (threshold $\geq 0.30$).
T4: RI = 1.000.
\textbf{Grade~A+.}

%----------------------------------------------------------------------
\section{Validation Results Summary}

\begin{table}[h]
\centering
\caption{Summary of all species validation results.}
\begin{tabular}{lrllllll}
\toprule
Species & Neurons & T1 & T2 $r_s$ & T3 & T4 RI & 30\,\% ret. & Grade \\
\midrule
\textit{D.\ rerio}        & 7,316 & PASS & 0.81 & PASS & 1.000 & 100\,\% & \textbf{A+} \\
\textit{C.\ elegans}      &   302 & PASS & 0.72 & PASS & 0.960 &  96\,\% & \textbf{A+} \\
\textit{D.\ melanogaster} & 2,145 & PASS & 0.68 & PASS & 0.946 &  89\,\% & \textbf{A+} \\
\textit{X.\ laevis}       &   150 & PASS & 1.00 & PASS & 1.000 & 100\,\% & \textbf{A+} \\
\bottomrule
\end{tabular}
\end{table}

%----------------------------------------------------------------------
\section{Usage}

\subsection{Run a full validation}

\begin{lstlisting}
from vzlab.validation.report import run_full_validation

report = run_full_validation("danio_rerio")
print(report.overall_grade())   # "A+"
print(report.robustness_index())  # 1.0
report.print_summary()
\end{lstlisting}

\subsection{CLI}

\begin{lstlisting}[language=bash]
vzlab run --species danio_rerio \
          --protocol predator_avoidance \
          --ablate amygdala \
          --episodes 20

vzlab sweep --species danio_rerio \
            --protocol exploration \
            --param efe.beta_epistemic \
            --values 0.0,1.0,2.0,3.0
\end{lstlisting}

\subsection{Register a new species}

\begin{lstlisting}
from vzlab.core.registry import register

register(
    name="mus_musculus",
    brain_cls=MouseBrain,
    env_cls=OpenFieldArena,
    connectome_path="data/mouse_connectome.json",
)
# CLI and ExperimentRunner discover it automatically.
\end{lstlisting}

%----------------------------------------------------------------------
\section{Biological Grounding of Design Decisions}

\paragraph{Sensory dropout for fault tolerance.}
Biological neural systems tolerate single-neuron loss, axonal
conduction failure, and synaptic unreliability because information is
represented redundantly across populations.
Tier~4 tests this property directly: an avatar that fails under 30\,\%
sensory dropout has learned brittle, non-redundant representations
that would not survive routine biological noise.

\paragraph{STDP synaptic dropout.}
Cortical and cerebellar synapses have vesicle release probabilities
of $\sim$0.1--0.3, meaning 70--90\,\% of spikes fail to transmit.
Applying Bernoulli dropout ($p = 0.10$) to the STDP weight update
prevents co-adaptation of co-active synapses and forces the learning
rule to distribute credit across many redundant pathways.

\paragraph{Spearman rank, not absolute values.}
Atlas activity vectors from calcium imaging or electrophysiology
contain unknown amplitude scaling factors.
Rank correlation tests the \emph{ordering} of regional activity
(which region is most vs.\ least active), not absolute firing rates.
This makes T2 robust to gain differences between the model and
measurement modality.

\paragraph{Lesion as ground truth.}
T3 replicates specific ablation experiments from the literature
rather than fitting to summary statistics.
A model that passes T3 has reproduced a causal mechanistic result,
not merely matched a correlation.

%----------------------------------------------------------------------
\section*{References}

\noindent
Roberts~A., Soffe~S.R., Wolf~E.S., Yoshida~M., Zhao~F.Y. (1998).
Central circuits controlling locomotion in young frog tadpoles.
\textit{Ann.\ NY Acad.\ Sci.}, 860, 19--34.

\noindent
Li~W.C., Soffe~S.R., Wolf~E., Roberts~A. (2006).
Persistent responses to brief stimuli: feedback excitation among
brainstem neurons.
\textit{J.\ Neurosci.}, 26(15), 4026--4035.

\noindent
Lin~A.C., Bygrave~A.M., de Calignon~A., Lee~T., Miesenb\"{o}ck~G.
(2014). Sparse, decorrelated odor coding in the mushroom body
enhances learned odor discrimination.
\textit{Nature Neurosci.}, 17, 559--568.

\noindent
White~J.G., Southgate~E., Thomson~J.N., Brenner~S. (1986).
The structure of the nervous system of the nematode
\textit{Caenorhabditis elegans}.
\textit{Phil.\ Trans.\ R.\ Soc.\ B}, 314, 1--340.

\noindent
Dorkenwald~S. \textit{et~al.} (2024). FlyWire: online community for
whole-brain connectomics.
\textit{Nature Methods}, 19, 119--128.

\end{document}
